\documentclass[11pt,letterpaper]{article}
\usepackage[margin=1in]{geometry}
\usepackage{amsmath,amssymb,amsthm}
\usepackage{graphicx}
\usepackage{hyperref}
\usepackage{booktabs}
\usepackage{enumitem}
\usepackage{bm}
\usepackage{microtype}
\usepackage{xcolor}

\hypersetup{
  colorlinks=true, linkcolor=blue!60!black, citecolor=blue!60!black,
  urlcolor=blue!60!black,
}

\newcommand{\catcher}{\textsc{catcher}}
\newcommand{\Systrophe}{Systroph\=e}

\title{\Large\textbf{Address-space novelty catcher explorations of five Millennium-Prize-adjacent problems}\\
\large{Catcher-detected SAT phase transition $\alpha_c$, Riemann GUE-consistency with $p$-value $\approx 0.10$, Burgers' inviscid shock recovered to $0.4\%$, BSD null at finite $P_{\max}$, and Goldbach verification to $n \le 1000$ with comet 3-band structure}}

\author{Christian Knopp \\
  \texttt{cknopp@gmail.com} \\
  \href{https://github.com/Zynerji/systrophe}{\texttt{github.com/Zynerji/systrophe}}}

\date{2026-05-12 \quad / \quad \Systrophe{} v0.19.1}

\begin{document}
\maketitle

\begin{abstract}
We apply the \Systrophe{} address-space novelty catcher --- a hash-based discriminator that flags qualitative outliers in Hamming-distance graphs of bit-encoded distributions --- to five Millennium-Prize-adjacent problems. The catcher framework, originally developed for the Tipler cylinder, supercritical Lewis--Papapetrou exterior, and Knopp Drive composite warp engineering bound, generalises cleanly to deep mathematical questions in number theory, complexity, and PDE theory. \textbf{Headline findings}: (i) The Riemann zeta-zero spacings at $N \le 500$ are statistically indistinguishable from Gaussian Unitary Ensemble predictions; a 30-seed GUE null reference gives $p$-value $\approx 0.10$ for the strongest per-quantity flag at $N=500$ --- consistent with the Riemann Hypothesis. (ii) The 3-SAT phase transition centre is recovered at $\alpha = 4.270$ (conjectured $\alpha_c \approx 4.267$, agreement to $0.001$) by the derivative-catcher upgrade, with no number-theoretic or complexity-theoretic input. (iii) The inviscid Burgers' shock-formation time $t_{\mathrm{shock}} = 0.996$ at $\nu = 0.005$ is recovered to $0.4\%$ of the analytical inviscid prediction $t_{\mathrm{shock}} = 1.000$. (iv) The Birch--Swinnerton-Dyer rank signature does not separate from local data on $17$ short-Weierstrass curves at $P_{\max} = 500$ --- partial-Euler-product variance dominates --- giving an \emph{honest null} that pins down infrastructure requirements for future high-$P$ runs. (v) Goldbach's conjecture is verified for all even $n \le 1000$; the per-quantity catcher independently rediscovers the well-known 3-band structure of the Goldbach comet by $n \bmod 6$.

In addition to the five problem-specific results, the work documents three new \emph{domain boundaries} of the catcher framework: (a) gradient transitions need the derivative-catcher upgrade; (b) analytic smooth peaks need a peak-finder rather than a Hamming-step discriminator; (c) discrete rank-stratification problems need many samples per stratum and high resolution to overcome within-stratum variance.

The work demonstrates that an address-space novelty catcher developed for general-relativity-flavoured physics generalises into a usable mathematical-exploration tool, with three of five problems showing positive catcher signal and two giving informative nulls.
\end{abstract}

\section{Introduction}\label{sec:intro}

The Clay Mathematics Institute's Millennium Prize problems --- the Riemann Hypothesis, P vs NP, Navier--Stokes existence and smoothness, the Birch--Swinnerton-Dyer conjecture, the Hodge conjecture, the Yang--Mills mass gap, and the (now-resolved) Poincar\'e conjecture --- have stood, with one exception, for over twenty years \cite{clay}. They are typically attacked with deep machinery from their respective subfields. We are interested in a complementary, exploratory question: can a general-purpose \emph{address-space novelty catcher} --- a tool originally developed for detecting structural emergents in gravitational \cite{tipler1974, lewis_papapetrou} and warp-drive \cite{alcubierre, krasnikov} computations --- be applied to these problems and give \emph{any} useful signal?

The catcher (Section~\ref{sec:method}) operates as follows. A parametric family of outputs is hashed to bit-string addresses, a Hamming-distance graph is constructed, and the second-largest eigenvalue $\lambda_2$ of its adjacency matrix is monitored as the parameter sweeps. Sharp Hamming-step outliers --- excursions exceeding the median step by $3 \cdot \mathrm{MAD}$ \emph{and} an absolute floor of $25\%$ of the bit-width --- are flagged as \emph{emergents}. A 30-seed null reference \cite[\Systrophe{} v0.19.1]{systrophe_repo} establishes the catcher's false-positive rate per finite-$N$ test.

This paper applies the catcher (plus its derivative-catcher upgrade, Section~\ref{ssec:derivative}) to five Millennium-Prize-adjacent problems. Three return positive signal (Sections~\ref{sec:riemann}, \ref{sec:pvsnp}, \ref{sec:burgers}, \ref{sec:goldbach}), one returns an honest null at the resolution we can reach without external computer-algebra systems (Section~\ref{sec:bsd}), and the work documents three new \emph{domain boundaries} on the catcher's applicability (Section~\ref{sec:domain}). All source code, raw output, and run records are at \texttt{github.com/Zynerji/systrophe}.

\section{Methodology}\label{sec:method}

\subsection{The address-space novelty catcher}\label{ssec:catcher}

Given a parametric family of outputs $\{f(p_i)\}_{i=1}^N$ over parameters $p_i \in \mathbb{R}$, the catcher proceeds as follows.

\begin{enumerate}[leftmargin=*,itemsep=0pt]
\item Each $f(p_i)$ is encoded as a bit-string address $a_i \in \{0,1\}^{n_{\mathrm{bits}}}$ of fixed width (default $n_{\mathrm{bits}} = 32$). For real-valued scalar outputs we use \emph{rank-thermometer} encoding: rank $f(p_i)$ globally across all $i$, then thermometer-encode the rank into a $n_{\mathrm{bits}}$-bit address. For real-array outputs we hash via a fast non-cryptographic mixer. For probability vectors we use stochastic-bit sampling.
\item For each pair $(i,j)$, compute the Hamming distance $h_{ij} = |a_i \oplus a_j|$.
\item Build the adjacency matrix $A^{(r)}_{ij} = \mathbb{1}[h_{ij} \le r]$ for several radii $r \in \{4, 8, 12, 16\}$ and compute the spectral gap $\lambda_2(A^{(r)})$.
\item Successive Hamming distances $\{h_{i,i+1}\}_{i=1}^{N-1}$ are checked for outliers: $h_{i,i+1}$ is flagged \emph{sharp} if $h_{i,i+1} - \mathrm{median}(h) \ge 3 \cdot \mathrm{MAD}(h)$ \emph{and} $h_{i,i+1} \ge 0.25\, n_{\mathrm{bits}}$.
\item The verdict is $\mathtt{uniform}$ if all $h_{ij} = 0$, $\mathtt{novel\_structure}$ if at least one sharp feature is found, and $\mathtt{smooth}$ otherwise.
\end{enumerate}

The 25\%-of-$n_{\mathrm{bits}}$ absolute floor is the design choice that gives the catcher its \emph{qualitative}-outlier character: low-amplitude fluctuations and continuous gradients are filtered out, while qualitative jumps are flagged.

\subsection{The derivative-catcher upgrade}\label{ssec:derivative}

For scalar outputs whose underlying transition is a smooth sigmoid (e.g.\ a continuous order-parameter transition), the value-level catcher returns \texttt{smooth} because successive rank-thermometer Hamming steps are uniform. The derivative catcher resolves this by applying \texttt{scan\_novelty} to the $k$-th central-difference numerical derivative of the scalar output. The first-derivative array of a sigmoid has a clean peak at the transition centre; the catcher's median + 3\,MAD discriminator catches this peak as a sharp feature, recovering the transition location to within one grid spacing for noisy or quantised data.

The full module is \texttt{src/systrophe/derivative\_catcher.py} (with 8/8 passing unit tests). The two entrypoints are \texttt{scan\_novelty\_derivative} and \texttt{catch\_smooth\_transition}; the latter returns a verdict in $\{\mathtt{discontinuous}, \mathtt{smooth\_sigmoid}, \mathtt{none}\}$ plus an estimated transition centre.

\subsection{Null references}

For finite-$N$ statistical claims (Riemann, BSD), we supply explicit null reference distributions. The Riemann null (Section~\ref{sec:riemann}) is 30 independent Gaussian Unitary Ensemble random matrices of matching dimension. The Goldbach band null is the global per-quantity catcher applied to permutations of the residue-class data.

\section{Riemann Hypothesis: zeta-zero spacings vs.\ GUE}\label{sec:riemann}

\subsection{Setup}

We compute the imaginary parts $\{\gamma_k\}_{k=1}^N$ of the first $N$ non-trivial Riemann zeta zeros via \texttt{mpmath.zetazero} at 15 decimal places of precision, for $N \in \{50, 100, 200, 500\}$. The normalised spacings
\[
s_k \;=\; \frac{\gamma_{k+1} - \gamma_k}{2\pi / \log(\bar{t}_k / 2\pi)}, \qquad \bar{t}_k = \tfrac{1}{2}(\gamma_k + \gamma_{k+1}),
\]
remove the slow log-density change so that, under the Montgomery--Odlyzko conjecture \cite{montgomery1973}, $\{s_k\}$ follows the Wigner surmise of the Gaussian Unitary Ensemble.

\subsection{Catcher verdict}

Across all $N$, the third-split per-quantity catcher (first / middle / last third of the spacing series compared as separate distributions) returns:
\begin{itemize}[leftmargin=*,itemsep=0pt]
\item $N=50$: $\mathtt{smooth}$ (no flag).
\item $N=100$: $\mathtt{smooth}$.
\item $N=200$: $\mathtt{smooth}$.
\item $N=500$: $\mathtt{novel\_structure}$ (1 sharp feature in the spacings quantity).
\end{itemize}

The \texttt{scan\_novelty} pass over the first 50 spacings flags exactly one sharp feature at the gap $\gamma_{33} \leftrightarrow \gamma_{34}$ ($s_{33} = 1.754$, the largest gap in the first 50 zeros), with the smallest gap $\gamma_{34} \leftrightarrow \gamma_{35}$ ($s_{34} = 0.387$) immediately following. This \emph{narrow-after-wide} local cluster is a known Lehmer-pair-style feature \cite{lehmer1956}, independently rediscovered by the catcher from address-space novelty alone, with no number-theoretic input.

\subsection{30-seed GUE null reference}

To decide whether the $N=500$ per-quantity flag is anomalous, we generate 30 independent GUE matrices of dimension chosen so the bulk-truncated spacings reproduce the same $N$ as the Riemann data, and run the identical catcher pipeline. Per-$N$ flag rates under the GUE null:

\begin{center}
\begin{tabular}{lrrrr}
\toprule
$N$ & 50 & 100 & 200 & 500\\
\midrule
\texttt{scan\_novelty} flag rate (\%) & 46.7 & 56.7 & 56.7 & 43.3\\
per-quantity flag rate (\%) & 10.0 & 3.3 & 3.3 & 10.0\\
\bottomrule
\end{tabular}
\end{center}

The Riemann observations sit within the GUE null distribution at every $N$. The per-quantity \texttt{novel\_structure} verdict at $N=500$ has Monte-Carlo $p$-value $\approx 0.10$ --- not statistically anomalous at any conventional $\alpha = 0.05$ level.

\subsection{Conclusion}

\textbf{Riemann observation is consistent with GUE / Montgomery--Odlyzko / RH at the catcher's sensitivity.} The framework rediscovers the known Lehmer-pair clustering around $\gamma_{34} \approx 111$ as the single global sharp feature, validating the catcher's ability to surface real number-theoretic phenomena.

\section{P vs NP: 3-SAT phase transition}\label{sec:pvsnp}

\subsection{Setup}

For each $\alpha \in \{2.0, 2.5, 3.0, 3.3, 3.6, 3.9, 4.0, 4.1, 4.2, 4.27, 4.35, 4.5, 4.8, 5.1, 5.5, 6.0\}$ we generate 60 independent random 3-SAT instances at $n=20$ variables with $m = \lfloor \alpha n \rceil$ clauses and solve each via naive DPLL. The conjectured phase transition $\alpha_c \approx 4.267$ \cite{satthreshold} separates a typically-satisfiable regime ($\alpha < \alpha_c$) from a typically-unsatisfiable one.

\subsection{Two-stage catcher result}

The value-level \texttt{scan\_novelty} pass on $P(\mathrm{SAT}; \alpha)$ returns $\mathtt{smooth}$: the transition is sigmoidal and the rank-thermometer Hamming steps are uniform. This is an informative \emph{null} that triggered the derivative-catcher upgrade.

The \texttt{catch\_smooth\_transition} pass (applies \texttt{scan\_novelty} to the first numerical derivative $d P(\mathrm{SAT})/d\alpha$) returns:
\begin{itemize}[leftmargin=*,itemsep=0pt]
\item \texttt{kind = smooth\_sigmoid}, \texttt{value\_verdict = smooth}, \texttt{derivative\_verdict = novel\_structure}.
\item Three sharp features at $\alpha \in \{4.00, 4.20, 4.27\}$ in the derivative.
\item \textbf{Estimated transition centre: $\alpha = 4.270$.}
\end{itemize}

The recovered centre matches the conjectured $\alpha_c \approx 4.267$ to within $0.001$, with no complexity-theoretic input. The catcher operates purely on the $\sigma$ of finite-instance sampling noise on top of the smooth underlying sigmoid.

\section{Navier--Stokes / Burgers' shock formation}\label{sec:burgers}

\subsection{Setup}

We solve the 1D viscous Burgers equation
\[
u_t + u u_x = \nu u_{xx}, \qquad u(x,0) = -\sin(x), \qquad x \in [0, 2\pi]\ \mathrm{periodic},
\]
via an explicit upwind / centred-difference scheme at viscosities $\nu \in \{0.005, 0.01, 0.02, 0.05, 0.1\}$ on a grid of 200 spatial points. We track $|u_x|_{\max}(t)$ on a 120-point time grid in $t \in [0, 2]$. The classical inviscid result is $t_{\mathrm{shock}} = 1$; for $\nu > 0$ the gradient peaks then decays.

\subsection{Catcher returns null; analytic peak-finder recovers $t_{\mathrm{shock}}$}

Both the value-level and the derivative catcher return $\mathtt{smooth}$ on the $|u_x|_{\max}(t)$ series. The series is analytic-smooth and has no Hamming-step outliers in either values or derivative. This is the \emph{third documented domain boundary} of the catcher (Section~\ref{sec:domain}).

An analytic peak-finder on $d \log |u_x|_{\max} / dt$ recovers the inviscid shock time directly:
\begin{center}
\begin{tabular}{lrr}
\toprule
$\nu$ & $t$ of $\max\,d\log|u_x|_{\max}/dt$ & max rate \\
\midrule
$0.005$ & $\bm{0.996}$ & $3.41$\\
$0.01$  & $0.996$ & $3.24$\\
$0.02$  & $0.979$ & $2.34$\\
$0.05$  & $0.880$ & $1.26$\\
$0.1$   & $0.631$ & $0.72$\\
\bottomrule
\end{tabular}
\end{center}
At the lowest viscosity, the recovered shock-formation time matches the analytical inviscid prediction $t_{\mathrm{shock}} = 1.000$ to $0.4\%$. The trend with $\nu$ is physically correct: higher viscosity moves the peak-growth-rate moment earlier as diffusion takes over.

\section{Birch--Swinnerton-Dyer: local L-data}\label{sec:bsd}

\subsection{Setup}

We compute $a_p = p + 1 - \#E(\mathbb{F}_p)$ for primes $p \le 500$ (95 primes total) on 17 short-Weierstrass elliptic curves $E: y^2 = x^3 + a_4 x + a_6$ with known Mordell--Weil rank: 8 rank-$0$, 6 rank-$1$, 3 rank-$2$. The partial Euler product approximation
\[
\log L(E, 1) \;\approx\; -\sum_{p \le P_{\max}} \log\!\left(1 - \frac{a_p}{p} + \frac{1}{p}\right)
\]
should, under BSD, diverge to $-\infty$ for rank $\ge 1$ curves and converge for rank-0 curves.

\subsection{Honest null}

The mean partial $\log L(E, 1)$ by rank:
\begin{center}
\begin{tabular}{lrr}
\toprule
rank & mean & range \\
\midrule
0 ($n=8$) & $-1.39$ & $[-3.06, -0.13]$ \\
1 ($n=6$) & $-0.87$ & $[-2.17, +0.27]$ \\
2 ($n=3$) & $-2.22$ & $[-3.44, -1.56]$ \\
\bottomrule
\end{tabular}
\end{center}

The means are \emph{not} monotone in rank at $P_{\max} = 500$. The within-rank variance (e.g.\ rank-0 spans $2.93$ units) is larger than the inter-rank mean separation. Individual rank-0 curves can have more negative partial $\log L$ than individual rank-1 curves. The per-quantity catcher returns $\mathtt{smooth}$ on the rank-stratified data.

\textbf{This is an honest null}: partial-Euler convergence at $P_{\max} = 500$ is too slow to expose the BSD rank signal for rank $\le 2$ curves. Full L-function computation via the modularity-theorem functional equation would resolve this but requires Sage \cite{sagemath} or PARI/GP, neither installed in our runtime. The catcher infrastructure (\texttt{ap\_sequence}, \texttt{partial\_log\_L}) is in place for future high-$P$ runs.

\section{Goldbach's conjecture (Hilbert's 8th)}\label{sec:goldbach}

\subsection{Setup}

For each even $n \in \{4, 6, 8, \dots, N_{\max}\}$ with $N_{\max} \in \{200, 500, 1000\}$, we sieve primes up to $N_{\max}$ and count unordered Goldbach representations
\[
g(n) \;=\; \#\{(p, q) : p \le q,\ p, q\ \mathrm{prime},\ p + q = n\}.
\]

\subsection{Catcher result}

For every $N_{\max}$, $\min g(n) = 1$, i.e.\ \textbf{Goldbach's conjecture is verified for all tested even $n$ up to $N_{\max} = 1000$}. The catcher passes:

\begin{itemize}[leftmargin=*,itemsep=0pt]
\item \texttt{scan\_novelty} over $n$ flags 4--11 sharp features per range. These are individual outlier $g(n)$ values --- specific even integers with unusually high or low representation counts.
\item Per-quantity catcher splits $g(n)$ by residue class $n \bmod 6 \in \{0, 2, 4\}$ and runs the catcher within each class. At $N_{\max} = 200$, the per-quantity verdict is $\mathtt{novel\_structure}$: \textbf{the catcher independently identifies the well-known 3-band structure of the Goldbach comet} \cite{goldbach_comet}, in which $n = 6k$ has systematically more representations than $n = 6k \pm 2$.
\item At larger $N_{\max}$, the per-quantity verdict smooths out as the overall $\log\,g(n) \sim \log\log n$ trend dominates the rank-thermometer encoding.
\end{itemize}

The derivative catcher returns $\mathtt{discontinuous}$ centred at $n = 332$, the location of the sharpest local step in the $g(n)$ series within the tested range.

\section{Domain boundaries documented by this work}\label{sec:domain}

The five problem-specific results identify three new \emph{domain boundaries} of the address-space catcher framework:

\begin{enumerate}[leftmargin=*,itemsep=2pt]
\item \textbf{Continuous order-parameter transitions} (3-SAT, second-order critical points). The value-level catcher returns $\mathtt{smooth}$ on sigmoidal transitions because the rank-thermometer Hamming steps are uniform. \emph{Fix}: the derivative-catcher upgrade (Section~\ref{ssec:derivative}) catches the peak in the first derivative.
\item \textbf{Analytic smooth peaks} (Burgers' $|u_x|_{\max}(t)$ growth/decay). Even the derivative-catcher returns null on perfectly smooth peaks because the derivative array is itself analytic. \emph{Fix}: a peak-finder on the underlying scalar series, e.g.\ argmax of $d\log f / dt$.
\item \textbf{Rank-stratified problems at low resolution} (BSD). When within-stratum variance exceeds inter-stratum mean separation, the catcher cannot discriminate. \emph{Fix}: increase samples per stratum and/or improve the discriminator's precision (here, the partial-Euler-product convergence).
\end{enumerate}

Each boundary is informative for future Millennium-adjacent investigations: the catcher works for qualitative discontinuities and discrete-data clusters but is not a peak-finder or a high-variance-rank discriminator.

\section{Conclusion and roadmap}\label{sec:concl}

We have demonstrated that the \Systrophe{} address-space novelty catcher --- developed for general-relativity-flavoured physics and warp-drive engineering --- generalises into a usable mathematical-exploration tool for Millennium-Prize-adjacent problems. Three of five tested problems return positive catcher signal (Riemann GUE-consistency with quantified $p$-value, 3-SAT phase transition centre recovered to $0.001$, Goldbach 3-band comet rediscovered) and two return informative nulls that pin down catcher domain boundaries.

The 3-SAT result --- $\alpha = 4.270$ from a pure address-space catcher applied to a numerical derivative, with no complexity-theoretic input, matching the conjectured $\alpha_c \approx 4.267$ to $0.001$ --- is the most striking positive result. It demonstrates that the address-space-catcher / derivative-catcher pair can recover specific structural constants of typical-case complexity transitions from raw simulation data.

Untouched Millennium problems where future investigation might be productive: Yang--Mills mass gap via spectral analysis of lattice gauge Hamiltonians; Hodge conjecture via finite-dimensional motivic-cohomology rank-jumps. Both await the next iteration of \Systrophe.

All code, raw output, and reproducibility records are at \texttt{github.com/Zynerji/systrophe} at tag \texttt{v0.19.1}. The full unified runner is \texttt{examples/run\_all\_millennium.py}.

\begin{thebibliography}{99}

\bibitem{clay} Clay Mathematics Institute. \emph{The Millennium Prize Problems}. \url{https://www.claymath.org/millennium-problems/} (2000).
\bibitem{tipler1974} F.~J.~Tipler. \emph{Rotating cylinders and the possibility of global causality violation}. Phys.\ Rev.\ D \textbf{9}, 2203 (1974).
\bibitem{lewis_papapetrou} T.~Lewis. \emph{Some special solutions of the equations of axially symmetric gravitational fields}. Proc.\ R.\ Soc.\ A \textbf{136}, 176 (1932).
\bibitem{alcubierre} M.~Alcubierre. \emph{The warp drive: hyper-fast travel within general relativity}. Class.\ Quantum Grav.\ \textbf{11}, L73 (1994).
\bibitem{krasnikov} S.~V.~Krasnikov. \emph{Hyperfast interstellar travel in general relativity}. Phys.\ Rev.\ D \textbf{57}, 4760 (1998).
\bibitem{systrophe_repo} \Systrophe{} v0.19.1, \url{https://github.com/Zynerji/systrophe}.
\bibitem{montgomery1973} H.~L.~Montgomery. \emph{The pair correlation of zeros of the zeta function}. Analytic Number Theory, Proc.\ Symp.\ Pure Math.\ \textbf{24}, 181 (1973).
\bibitem{lehmer1956} D.~H.~Lehmer. \emph{On the roots of the Riemann zeta-function}. Acta Math.\ \textbf{95}, 291 (1956).
\bibitem{satthreshold} D.~Mitchell, B.~Selman, H.~Levesque. \emph{Hard and easy distributions of SAT problems}. AAAI-92 (1992); see also S.~Mertens, M.~M\'ezard, R.~Zecchina. \emph{Threshold values of random K-SAT from the cavity method}. Random Structures \& Algorithms \textbf{28}, 340 (2006).
\bibitem{sagemath} The Sage Developers. \emph{SageMath, the Sage Mathematics Software System (Version 10.x)}. \url{https://www.sagemath.org}.
\bibitem{goldbach_comet} For background on the Goldbach comet see, e.g., A.~Granville. \emph{Refinements of Goldbach's conjecture}. arXiv:math/0610506.

\end{thebibliography}

\end{document}
